% Part: history
% Chapter: biographies
% Section: ernst-zermelo

\documentclass[../../../include/open-logic-section]{subfiles}

\begin{document}

\olfileid{his}{bio}{zer}

\olsection{અર્ન્સ્ટ ઝર્મેલો}

\olphoto{zermelo-ernst}{અર્ન્સ્ટ ઝર્મેલો}

અર્ન્સ્ટ ઝર્મેલોનો જન્મ 27 જુલાઈ
1871ના રોજ જર્મનીના બર્લિનમાં
થયો હતો. તેમને પાંચ બહેનો
હતી, પરંતુ કુટુંબમાં નબળા
સ્વાસ્થ્યને કારણે તેમાંથી
માત્ર ત્રણ જ પુખ્ત વયે
પહોંચી. ઝર્મેલો નાના હતા
ત્યારે જ તેમનાં માતાપિતાનું
પણ અવસાન થયું. તેઓ સત્તર
વર્ષના હતા ત્યારે તેઓ અને
તેમની બહેનો અનાથ થઈ
ગયાં હતાં. ઝર્મેલોને કલાઓમાં,
ખાસ કરીને કવિતામાં,
ઊંડો રસ હતો. તેઓ
તીક્ષ્ણબુદ્ધિ, વિનોદી અને
ટીકાત્મક સ્વભાવના
તરીકે જાણીતા હતા.
તેમની સૌથી જાણીતી
ગાણિતિક સિદ્ધિઓમાં
1904માં પસંદગીની
સ્વયંસિદ્ધિની રજૂઆત
અને 1908માં ગણસિદ્ધાંતનું
સ્વયંસિદ્ધીકરણ સામેલ છે.

યુનિવર્સિટીમાં ઝર્મેલોની
રુચિઓ વિવિધ હતી. તેમણે
ભૌતિકશાસ્ત્ર, ગણિત અને
તત્વજ્ઞાનના અભ્યાસક્રમો
લીધા. હર્માન શ્વાર્ઝના
માર્ગદર્શન હેઠળ તેમણે
1894માં બર્લિન
યુનિવર્સિટીમાં
\emph{Investigations in
  the Calculus of Variations}
નામનો સંશોધનનિબંધ
પૂરો કર્યો. 1897માં
તેમણે ગોટિંગેન
યુનિવર્સિટીમાં વધુ
અભ્યાસ કરવાનું નક્કી
કર્યું; ત્યાં ડેવિડ
હિલ્બર્ટના પાયાકીય
કાર્યનો તેમના પર
ઘણો પ્રભાવ પડ્યો.
1899માં તેઓ પ્રોફેસરપદ
માટે લાયક બન્યા,
પરંતુ અગિયાર વર્ષ
પછી જ તેમને એ પદ
મળ્યું—કદાચ તેમના
વિચિત્ર વર્તન અને
“ઉતાવળભરી અધીરાઈ”ને
કારણે.

આખરે 1910માં ઝર્મેલોને
ઝ્યુરિખ યુનિવર્સિટીમાં
પગારવાળું પ્રોફેસરપદ
મળ્યું, પરંતુ ક્ષયરોગને
કારણે 1916માં નિવૃત્ત
થવાની ફરજ પડી.
સ્વસ્થ થયા પછી 1921માં
તેમને ફ્રાઇબર્ગ
યુનિવર્સિટીમાં માનદ
પ્રોફેસરપદ મળ્યું.
આ સમય દરમિયાન તેમણે
ગણિતના પાયા પર કામ
કર્યું. થોરાલ્ફ સ્કોલેમ
અને કુર્ત ગેડલનાં
કાર્યથી તેઓ ખીજાયા
અને પોતાના લેખોમાં
તેમના અભિગમોની
જાહેરમાં ટીકા કરી.
1935માં પોતાની
અલોકપ્રિયતા અને
જર્મનીમાં હિટલરના
સત્તારોહણના વિરોધને
કારણે તેમને ફ્રાઇબર્ગના
પદેથી બરતરફ કરાયા.
 
ઝર્મેલોના જીવનનાં
પાછલાં વર્ષો એકલતામાં
વીત્યાં. 1935ની
બરતરફી પછી તેમણે
ગણિત છોડી દીધું.
તેઓ ગ્રામ્ય વિસ્તારમાં
રહેવા ગયા અને સાદગીથી
જીવ્યા. 1944માં તેમણે
લગ્ન કર્યાં. દૃષ્ટિ
ઓછી થતી જતાં તેઓ
પત્ની પર સંપૂર્ણપણે
નિર્ભર બન્યા. 1951
સુધીમાં તેમની દૃષ્ટિ
સંપૂર્ણ જતી રહી.
21 મે 1953ના રોજ
જર્મનીના ગુન્ટરસ્ટાલમાં
તેમનું અવસાન થયું.

\begin{reading}
ઝર્મેલોના સંપૂર્ણ
જીવનચરિત્ર માટે
\citet{Ebbinghaus2015}
જુઓ. તેમના 1904
અને 1908ના પાયાના
લેખો મૂળ જર્મન
ભાષામાં વાંચવા
ઉપલબ્ધ છે
\citep{Zermelo1904,Zermelo1908}.
ભૌતિકશાસ્ત્ર અંગેના
તેમના લખાણો સહિત
એકત્રિત કૃતિઓના
અંગ્રેજી અનુવાદો
\citep{Ebbinghaus2010,Ebbinghaus2013}માં
ઉપલબ્ધ છે.
\end{reading}

\end{document}
